\documentclass[10pt,a4paper]{amsart}
\usepackage[margin=19mm]{geometry}
\usepackage{amsmath,amssymb,mathtools}
\usepackage[scr=rsfs,cal=euler]{mathalfa}
\usepackage{xfrac}
\usepackage[unicode,hidelinks,bookmarks=false]{hyperref}
\newcommand{\ie}{i.e.\ }
\newcommand{\cf}{cf.\ }
\newcommand{\resp}{respectively\ }
\newcommand{\Subsection}{\subsection}
\providecommand{\nobreakdash}{\nobreak\hbox{-}}
\setlength{\emergencystretch}{2em}
\multlinegap0pt
\usepackage{amssymb}
\usepackage{mathtools}
\usepackage{bm}
\usepackage[shortlabels]{enumitem}
\newtheorem{corollary}{Corollary}
\newcommand{\HH}{\mathcal{H}}
\newcommand{\Ric}{\operatorname{Ric}}
\newcommand{\tg}{\widetilde{g}}
\title{Hineva Inequality for Submanifolds of Real Space Forms with Semi-Symmetric Non-Metric Connection}
\author{M. S. Lone, M. A. Lone}
\date{Source: arXiv:2606.05193v2 (CC BY 4.0)}
\begin{document}
\maketitle

\noindent\emph{Source and licence.} Hineva Inequality for Submanifolds of Real Space Forms with Semi-Symmetric Non-Metric Connection. Version arXiv:2606.05193v2 (CC BY 4.0), distributed by arXiv under the Creative Commons Attribution 4.0 International licence, http://creativecommons.org/licenses/by/4.0/, as recorded in the arXiv licence metadata for this exact version and shown on the abstract page https://arxiv.org/abs/2606.05193v2. Full licence notice: \texttt{licenses/arxiv-2606.05193-CC-BY-4.0.txt}.\par\smallskip

\noindent\textit{Editorial scope.} Counted unit: the article's corollary cor:Pnor (source0.tex lines 275--277). The article's complete original proof of it is delivered with it (source0.tex lines 278--280, one \texttt{proof} environment of 3 lines). No mathematics is written, repaired or re-derived: the statement and the proof are the article's own lines, in the article's own notation, and no retained line is substituted or re-flowed.  Retained uncounted as the source-local context the proof needs: lines 207--209.  Nothing was omitted from any retained span, so the package carries no omission record.  No retained line contains a citation, so the wrapper carries no bibliography and no bibliography entry.  No retained line contains a figure environment, an image-inclusion command or drawing code, so no image is delivered and the package declares no asset dependency.  Editorial formatting only, in addition: an AMS \texttt{amsart} preamble that restates the article's own declarations and macros used by the retained lines, namely the environments \texttt{corollary} on separate counters, together with the standard packages those lines need.  The retained environments print in this document as Corollary 1 in order of appearance, whereas the article prints the same items with the numbering of its own full document.  The complete native source of the article as distributed in the arXiv e-print for this exact version is bundled unchanged as \texttt{sources/202098/source0.tex}.
\begin{equation}\label{eq:mainineq}
\Ric(X)\;\ge\;\Theta(X)-\lambda+s(X,X)-n\,\phi(H)+\tg\bigl(P^\perp,\sigma(X,X)\bigr)+\HH(H,\sigma),
\end{equation}

\begin{corollary}[$P$ normal]\label{cor:Pnor}
If $P^\top=0$, then $\Ric(X)\ge\Theta(X)-\lambda+s(X,X)-n\,\tg(P,H)+\tg(P,\sigma(X,X))+\HH(H,\sigma)$.
\end{corollary}

\begin{proof}
Since $P^\top=0$ we have $P=P^\perp$, so $\phi(H)=\tg(P^\perp,H)=\tg(P,H)$ and $\tg(P^\perp,\sigma(X,X))=\tg(P,\sigma(X,X))$; substituting into \eqref{eq:mainineq} gives the result.
\end{proof}

\end{document}
